\documentclass[11pt]{amsart}
\usepackage[a4paper, margin=2cm]{geometry}
\usepackage{amsmath,amssymb,amsthm,mathtools}
\usepackage{booktabs}
\usepackage{enumitem}
\usepackage[hidelinks]{hyperref}
\frenchspacing

\newtheorem{theorem}{Theorem}[section]
\newtheorem{lemma}[theorem]{Lemma}
\newtheorem{proposition}[theorem]{Proposition}
\newtheorem{corollary}[theorem]{Corollary}
\theoremstyle{definition}
\newtheorem{remark}[theorem]{Remark}
\numberwithin{equation}{section}

\DeclareMathOperator{\Tr}{Tr}
\DeclareMathOperator{\Var}{Var}
\DeclareMathOperator{\N}{N}
\DeclareMathOperator{\Cl}{Cl}
\newcommand{\OK}{\mathcal{O}_K}
\newcommand{\OKd}{\mathcal{O}_K^{\vee}}
\newcommand{\Dif}{\mathfrak{D}}
\newcommand{\Q}{\mathbb{Q}}
\newcommand{\Z}{\mathbb{Z}}
\newcommand{\R}{\mathbb{R}}
\newcommand{\LamK}{\Lambda_K}
\newcommand{\one}{\mathbf{1}}
\newcommand{\file}[1]{\texttt{#1}}

\title{Siegel's formula by ideal classes and the lifting problem}
\date{Draft notes, November 2026. Code and data: \file{higher-degree/} in \url{https://github.com/pav10/2O}}

\begin{document}

\begin{abstract}
Kala and Yatsyna bounded the discriminant of a totally real field of degree $2$, $3$, $4$, $5$ or $7$
that admits a universal quadratic form with rational integer coefficients, under the assumption
that the codifferent is a principal ideal.  Kala and Melistas extended this to every degree up to $43$
when, in addition, the minimal trace of a totally positive element of the codifferent is as large as
Siegel's formula allows.  We remove the principal codifferent assumption from both results, with
explicit bounds.  The method splits Siegel's formula for $\zeta_K(-1)$ into one identity for every
ideal class and compares the class of the different with the classes of its divisors.  In degrees
$2$ and $3$ the bound is small enough for a complete check, which gives a short proof that
$\Q(\sqrt5)$ and $\Q(\zeta_7+\zeta_7^{-1})$ are the only real quadratic and cubic fields with such a
form.
\end{abstract}

\maketitle

\section{Introduction}

Let $K$ be a totally real number field of degree $d$, with ring of integers $\OK$, codifferent
$\OKd=\Dif^{-1}$ and discriminant $\Delta_K$.  A \emph{universal $\Z$-form} over $K$ is a positive
definite quadratic form with coefficients in $\Z$ that represents every totally positive element of
$\OK$.  The lifting problem asks for the fields that admit one.  Kala and Yatsyna conjectured that
these are $\Q$, $\Q(\sqrt5)$ and $\Q(\zeta_7+\zeta_7^{-1})$ \cite{KY1}.  The conjecture is known in
degrees at most $5$ \cite{KY2} and, by the computer-assisted notes in this folder, in degrees $6$,
$7$ and $8$.

For $\ell\ge1$ let
\[
  T_\ell=\{\nu\in\OKd : \nu\succ0,\ \Tr\nu=\ell\},\qquad m(K)=\min\{\ell : T_\ell\ne\emptyset\},
\]
and let $r_d=\dim M_{2d}(\mathrm{SL}_2(\Z))$.  Siegel's formula \cite{Si} writes $\zeta_K(-1)$ as a
linear combination of the sums $\sum_{\nu\in T_\ell}\sigma(\nu\Dif)$, $\ell\le r_d$, where
$\sigma(\mathfrak a)=\sum_{\mathfrak b\mid\mathfrak a}\N\mathfrak b$.  Kala and Yatsyna
\cite[Theorem~1.2]{KY1} observed that a universal $\Z$-form forces the totally positive elements of
$\OK$ that minimize $\Tr(\delta\,\cdot)$ to be units (via Kitaoka's theory of lattices of E-type), and
deduced that if $d\in\{2,3,4,5,7\}$ (exactly the degrees with $r_d=1$) and $\OKd$ is principal, then
$|\Delta_K|$ is bounded.  Kala and Melistas \cite[Theorem~2.4]{KM} proved the analogous bound for
$d\le43$ when $\OKd$ is principal and $m(K)=r_d$.  In both results principality makes $\nu\Dif$
principal for $\nu\in T_{m(K)}$, so that the minimizing property gives $\sigma(\nu\Dif)=1$ and bounds
$\#T_{m(K)}$ by a number of units.  Kim and Lee \cite{KL} solved the cubic case without this assumption by a
different method, and Kala and Melistas \cite[\S3]{KM} proposed to replace $\OKd$ by a sum of few
principal pieces.

We will keep $\nu\Dif$ non-principal and split Siegel's formula instead.  The universal $\Z$-form
enters only through the following property of the trace-minimal elements of the codifferent:
\begin{equation}\label{eq:condP}
  \text{for every }\nu\in T_{m(K)},\ \text{the only principal ideal dividing }\nu\Dif\text{ is }\OK.
\end{equation}
Lemma~\ref{lem:etype} shows that \eqref{eq:condP} holds when $d\le43$ and $K$ admits a universal
$\Z$-form.

\begin{theorem}\label{thm:main}
Let $K$ be a totally real field of degree $d\ge2$ with $m(K)=r_d$ that satisfies \eqref{eq:condP}.
Let $A_d$, $c_d$, $G_d$ be the constants of \eqref{eq:Ad}, Lemma~\textup{\ref{lem:count}} and
Lemma~\textup{\ref{lem:small}}.  Then:
\begin{enumerate}[label=\textup{(\arabic*)}]
\item\label{it:class} every integral ideal in the class of $\Dif$ has norm at least
  $(|\Delta_K|/(A_dc_d))^{1/2}$;
\item\label{it:disc} $|\Delta_K|\le(A_dc_dG_d^2)^{d-1}$;
\item\label{it:reg} the regulator satisfies $R_K\le\gamma_d^d A_dc_d/2^{d-2}$, where $\gamma_d$ is
  Hermite's constant.
\end{enumerate}
\end{theorem}

The condition $m(K)=r_d$ holds for every field of degree $2$, $3$, $4$, $5$ and $7$, since $r_d=1$
there and $m(K)\le r_d$ by Siegel's formula.  For $d=2$ and $d=3$ the bound in
\ref{it:disc} is $13$ and $3621$, and a complete check of these fields gives the following.

\begin{corollary}\label{cor:cubic}
The only real quadratic field that satisfies \eqref{eq:condP} is $\Q(\sqrt5)$, and the only totally
real cubic field that satisfies it is $\Q(\zeta_7+\zeta_7^{-1})$.  In particular these are the only real quadratic and
cubic fields with a universal $\Z$-form.
\end{corollary}

The second sentence was proved by Kala and Yatsyna in the quadratic case \cite[Theorem~1.1]{KY1}
and by Kim and Lee in the cubic case \cite[Theorem~1.3]{KL}.

The proof of Theorem~\ref{thm:main} will use three ingredients.  The first is Siegel's formula split
by ideal classes (Proposition~\ref{prop:siegel}).  It shows that \eqref{eq:condP} is equivalent to
an identity between $\#T_{r_d}$ and the partial zeta value $\zeta(2;[\Dif])$ of the class of the
different (Proposition~\ref{prop:identity}).  The second is a lattice-point bound
$\#T_{r_d}\ll|\Delta_K|^{1/2}$ (Lemma~\ref{lem:count}).  Together they say that $[\Dif]$ contains no
ideal of norm much smaller than $|\Delta_K|^{1/2}$.  The third ingredient contradicts this for large
$|\Delta_K|$: a shortest nonzero element $\xi$ of trace zero in $\OKd$ gives the integral ideal
$\xi\Dif\in[\Dif]$ of norm $O(|\Delta_K|^{(d-2)/(2d-2)})$ (Lemma~\ref{lem:small}).  So the lifting
problem in these degrees reduces to estimates for one partial zeta value.

The explicit bounds grow quickly with $d$ (Table~\ref{tab:bounds}), so only degrees $2$ and $3$ are
settled by a direct check.  In degrees $d\ge6$ with $m(K)<r_d$ the formula also contains the sums
over $T_\ell$ for $m(K)<\ell\le r_d$, which the minimizing property does not control.  This, not the
principality of the codifferent, is what remains.  We discuss it in Section~\ref{sec:higher}.

\section{Siegel's formula by ideal classes}\label{sec:siegel}

Throughout, $K$ is totally real of degree $d\ge2$, and we identify $K$ with its image in $\R^d$
under the real embeddings $\sigma_1,\dots,\sigma_d$.  Then the trace form is the standard inner
product, $\OKd$ is the dual lattice of $\OK$, and $\Dif=(\OKd)^{-1}$ has norm $|\Delta_K|$.  For
$\nu\in\OKd\setminus\{0\}$ the ideal $\nu\Dif$ is integral and lies in the class $[\Dif]$.  Let $U$ be
the unit group, $U^+$ its totally positive part, $\Cl$ the (wide) ideal class group, and for a class
$C$ put
\[
  \zeta(s;C)=\sum_{\mathfrak b\subseteq\OK,\ [\mathfrak b]=C}\N\mathfrak b^{-s}\quad(s>1),\qquad
  a_C(\nu)=\sum_{\mathfrak b\mid\nu\Dif,\ [\mathfrak b]=C}\N\mathfrak b .
\]
Let $b_1(2d),\dots,b_{r_d}(2d)$ be the rational numbers with $a_0(f)=\sum_{\ell=1}^{r_d}b_\ell(2d)a_\ell(f)$
for every $f=\sum_{n\ge0}a_n(f)q^n$ in $M_{2d}(\mathrm{SL}_2(\Z))$, see \cite{Si} or
\cite[Theorem~2.2]{KM}.

\begin{proposition}\label{prop:siegel}
For every class $C\in\Cl$,
\begin{equation}\label{eq:siegelC}
  \sum_{\ell=1}^{r_d}b_\ell(2d)\sum_{\nu\in T_\ell}a_C(\nu)
  =(-1)^d\,\frac{|\Delta_K|^{3/2}}{(4\pi^2)^d}\,\zeta\bigl(2;C^{-1}[\Dif]\bigr).
\end{equation}
\end{proposition}

Summing \eqref{eq:siegelC} over $C$ and using the functional equation
$\zeta_K(-1)=(-1)^d|\Delta_K|^{3/2}\zeta_K(2)/(2\pi^2)^d$ gives Siegel's formula
$\zeta_K(-1)=2^d\sum_\ell b_\ell(2d)\sum_{\nu\in T_\ell}\sigma(\nu\Dif)$.  We have not found the class
version in the literature.  We checked it numerically in degrees $3$, $4$ and $6$
(Section~\ref{sec:computations}).

\begin{proof}
For a fractional ideal $\mathfrak a$, an even integer $k\ge2$ and $z\in\mathcal H^d$ put
\[
  G_{\mathfrak a}(z)=\sum_{(c,e)\in(\mathfrak a\times\mathfrak a\setminus\{0\})/U}\N(cz+e)^{-k},
  \qquad \N(cz+e)=\prod_{i=1}^d\bigl(\sigma_i(c)z_i+\sigma_i(e)\bigr),
\]
where $U$ acts diagonally.  For $k=2$ the series is defined by Hecke's limiting process, exactly as
for Siegel's formula.  Since $d\ge2$, the limit is holomorphic and the computation below is unchanged.
The group $\mathrm{SL}_2(\OK)$ permutes $\mathfrak a\times\mathfrak a$, so $G_{\mathfrak a}$ is a
Hilbert modular form of parallel weight $k$ for $\mathrm{SL}_2(\OK)$, and its restriction to the
diagonal $z_1=\dots=z_d=\tau$ lies in $M_{kd}(\mathrm{SL}_2(\Z))$.

We compute the Fourier expansion.  The pairs with $c=0$ contribute
$\sum_{e\in(\mathfrak a\setminus0)/U}\N(e)^{-k}=\N\mathfrak a^{-k}\zeta(k;[\mathfrak a]^{-1})$.
For $c\ne0$ we use the Lipschitz formula (Poisson summation over $\OK$)
\[
  \sum_{\mu\in\OK}\N(z+\mu)^{-k}=\kappa_k|\Delta_K|^{-1/2}\sum_{\nu\in\OKd,\ \nu\succ0}\N(\nu)^{k-1}e^{2\pi i\Tr(\nu z)},
  \qquad \kappa_k=\Bigl(\frac{(-2\pi i)^k}{(k-1)!}\Bigr)^{d}.
\]
Writing $e=e_0+c\mu$ with $e_0$ in a set of representatives of $\mathfrak a/c\OK$ and $\mu\in\OK$, we
get
\[
  \sum_{e\in\mathfrak a}\N(cz+e)^{-k}=\N(c)^{-k}\kappa_k|\Delta_K|^{-1/2}\sum_{\nu\succ0}\N(\nu)^{k-1}e^{2\pi i\Tr(\nu z)}
  \sum_{e_0}e^{2\pi i\Tr(\nu e_0/c)} .
\]
The inner sum is $[\mathfrak a:c\OK]=|\N(c)|/\N\mathfrak a$ if $\nu/c\in(\mathfrak a\Dif)^{-1}$, and $0$
otherwise.  The condition says that the integral ideal $\mathfrak b=c\mathfrak a^{-1}$ divides
$\nu\Dif$, and as $c$ runs over $(\mathfrak a\setminus0)/U$ the ideal $\mathfrak b$ runs over the
integral ideals of the class $[\mathfrak a]^{-1}$, with $|\N(c)|=\N\mathfrak b\,\N\mathfrak a$.
Replacing $\mathfrak b$ by $\nu\Dif\mathfrak b^{-1}$ and using $\N(\nu\Dif)=\N(\nu)|\Delta_K|$, we
obtain
\[
  \N\mathfrak a^{k}\,G_{\mathfrak a}(\tau,\dots,\tau)
  =\zeta\bigl(k;[\mathfrak a]^{-1}\bigr)+\kappa_k|\Delta_K|^{1/2-k}\sum_{\ell\ge1}
  \Bigl(\sum_{\nu\in T_\ell}\ \sum_{\mathfrak b\mid\nu\Dif,\ [\mathfrak b]=[\mathfrak a\Dif]}\N\mathfrak b^{k-1}\Bigr)q^\ell .
\]
Take $k=2$, choose $\mathfrak a$ with $[\mathfrak a\Dif]=C$, so that $[\mathfrak a]^{-1}=C^{-1}[\Dif]$,
and apply $a_0=\sum_\ell b_\ell(2d)a_\ell$.  Since $\kappa_2=(-4\pi^2)^d$, this is \eqref{eq:siegelC}.
\end{proof}

Since $\zeta(2;C)>0$ for every class $C$, \eqref{eq:siegelC} shows that $T_\ell\ne\emptyset$ for some
$\ell\le r_d$, that is, $m(K)\le r_d$.

\section{The trace-minimal elements of the codifferent}\label{sec:etype}

\begin{lemma}\label{lem:etype}
Assume that $d\le43$ and that $K$ admits a universal $\Z$-form.  Then $K$ satisfies
\eqref{eq:condP}.
\end{lemma}

\begin{proof}
Let $\nu\in T_{m(K)}$ and let $\alpha\in\OK$ be such that $(\alpha)$ divides $\nu\Dif$.  By
\cite[Corollary~4.5]{KY1} there is a unit of every signature, so we may assume $\alpha\succ0$.  Then
$\delta=\nu/\alpha$ lies in $\OKd$ and is totally positive.  For every totally positive
$\beta\in\OK$ the element $\delta\beta$ lies in $T_\ell$ for some $\ell$, so
$\Tr(\delta\beta)\ge m(K)=\Tr(\delta\alpha)$.  By \cite[Lemma~4.6]{KY1}, $\alpha$ is a unit.
\end{proof}

Both results of \cite{KY1} that we used rest on Kitaoka's theorem that the trace form
$x\mapsto\Tr(\delta x^2)$ on $\OK$, which has rank $d\le43$, is of E-type \cite[Theorem~4.2]{KY1}.

\begin{proposition}\label{prop:identity}
Assume that $m(K)=r_d$ and put
\begin{equation}\label{eq:Ad}
  A_d=(4\pi^2)^d\,|b_{r_d}(2d)| ,\qquad\text{so that } A_d=(2d-1)!\,\zeta(2d)\ \text{if } r_d=1 .
\end{equation}
Then $\#T_{r_d}\le|\Delta_K|^{3/2}\zeta(2;[\Dif])/A_d$, with equality if and only if $K$ satisfies
\eqref{eq:condP}.  If $K$ satisfies \eqref{eq:condP}, then also
\[
  \sum_{\nu\in T_{r_d}}\N(\nu)=\frac{|\Delta_K|^{1/2}\,\zeta(2;[\OK])}{A_d}.
\]
\end{proposition}

\begin{proof}
Only $\ell=r_d$ occurs in \eqref{eq:siegelC}.  For $C=[\OK]$ we have $a_C(\nu)\ge1$ (the divisor
$\OK$), with equality for all $\nu\in T_{r_d}$ exactly when \eqref{eq:condP} holds.  Since
$T_{r_d}\ne\emptyset$ and $\zeta(2;[\Dif])>0$, \eqref{eq:siegelC} shows that $(-1)^db_{r_d}(2d)>0$,
and the first claim follows.  For $C=[\Dif]$ the map
$\mathfrak b\mapsto\nu\Dif\mathfrak b^{-1}$ is a bijection from the divisors of $\nu\Dif$ in $[\Dif]$
onto its principal divisors.  Under \eqref{eq:condP} the only such divisor is $\nu\Dif$ itself, so
$a_{[\Dif]}(\nu)=\N(\nu)|\Delta_K|$.  The formula for $A_d$ when $r_d=1$ follows from
$b_1(2d)=-B_{2d}/(4d)$ and Euler's formula for $\zeta(2d)$.
\end{proof}

If $\OKd$ is principal, then $\zeta(2;[\Dif])=\zeta(2;[\OK])\ge1$ and \eqref{eq:condP} forces
$T_{r_d}\subseteq\delta_0U^+$ for a generator $\delta_0$; this is the situation of
\cite{KY1,KM}.  In general $[\Dif]$ is the square of a class by Hecke's theorem (see \cite{Ar} and
\cite[Remark~1.2]{KL}), so $\OKd$ is principal whenever $\Cl$ is an elementary abelian $2$-group.

\section{Three estimates}\label{sec:estimates}

Let $H=\{y\in\R^d:\sum_iy_i=0\}$, let $\pi$ be the orthogonal projection onto $H$, and put
\[
  L_K=\OKd\cap H,\qquad \LamK=\pi(\OK).
\]
Here $|\pi(x)|^2=d\Var(x)$ with $\Var(x)=\Tr(x^2)/d-(\Tr(x)/d)^2$, and $\pi(x)=0$ exactly when
$x\in\Z$, so $\LamK$ is the variance lattice $\OK/\Z$ of the degree-$6$ and $7$ notes.  Since
$1$ extends to a $\Z$-basis of $\OK$, the trace maps $\OKd$ onto $\Z$, and its level sets are
$1/\sqrt d$ apart.  Hence $L_K$ is a lattice of rank $d-1$ in $H$ with covolume
$\sqrt d\,|\Delta_K|^{-1/2}$, and its dual lattice in $H$ is $\LamK$.

\begin{lemma}\label{lem:small}
Let
\[
  \eta_d=\max\Bigl\{\prod_{i=1}^d|y_i| : y\in H,\ |y|=1\Bigr\},\qquad
  G_d=\eta_d\bigl(\gamma_{d-1}d^{1/(d-1)}\bigr)^{d/2}.
\]
There is a nonzero $\xi\in L_K$ with $\N(\xi\Dif)\le G_d\,|\Delta_K|^{(d-2)/(2d-2)}$.  In particular
the class $[\Dif]$ contains an integral ideal of at most this norm.
\end{lemma}

\begin{proof}
By the definition of Hermite's constant there is a nonzero $\xi\in L_K$ with
$|\xi|^2\le\gamma_{d-1}(\sqrt d\,|\Delta_K|^{-1/2})^{2/(d-1)}$.  Then
$|\N(\xi)|\le\eta_d|\xi|^d$, and $\N(\xi\Dif)=|\N(\xi)|\,|\Delta_K|$.
\end{proof}

By Lagrange multipliers a maximizing $y$ takes two values, $d-k$ on $k$ coordinates and $-k$ on the
others (up to scaling), so $\eta_d=\max_{1\le k<d}(d-k)^kk^{d-k}/(dk(d-k))^{d/2}$.  For even $d$ this
is $d^{-d/2}$; for $d=3$ it is $1/(3\sqrt6)$.

\begin{lemma}\label{lem:count}
For $\ell\ge1$ we have $\#T_\ell\le c_{d,\ell}|\Delta_K|^{1/2}$, where
\[
  c_{d,\ell}=d^{-1/2}\Bigl(\frac{2\pi e\ell^2}{d}\Bigr)^{(d-1)/2}(1+\varepsilon_{d,\ell}),\qquad
  \varepsilon_{d,\ell}=a\int_{1/3}^\infty e^{-aV}\bigl((1+2\sqrt{3V})^{d-1}-1\bigr)\,dV,\quad a=2\pi^2\ell^2 .
\]
We write $c_d=c_{d,r_d}$.
\end{lemma}

\begin{proof}
Every $y\in\R^d$ with positive coordinates and $\sum_iy_i=\ell$ satisfies
$|y-\tfrac{\ell}{d}\one|^2=|y|^2-\ell^2/d\le\ell^2(d-1)/d$.  The elements of $\OKd$ of trace $\ell$ form
a translate $\tfrac{\ell}{d}\one+w_0+L_K$ with $w_0\in H$.  Hence, for every $t>0$,
\[
  \#T_\ell\le\sum_{w\in w_0+L_K}e^{\pi t(\ell^2(d-1)/d-|w|^2)}
  \le e^{\pi t\ell^2(d-1)/d}\,\frac{t^{-(d-1)/2}}{\sqrt d\,|\Delta_K|^{-1/2}}\sum_{\lambda\in\LamK}e^{-\pi|\lambda|^2/t},
\]
by Poisson summation over $L_K$, whose dual is $\LamK$.  We take $t=d/(2\pi\ell^2)$.  Then the first two
factors give $d^{-1/2}(2\pi e\ell^2/d)^{(d-1)/2}|\Delta_K|^{1/2}$, and the last sum is
$\sum_{x\in\OK/\Z}e^{-a\Var(x)}$.

We claim that $\Var(x)\ge1/3$ for $x\in\OK\setminus\Z$.  Let $x_1,\dots,x_n$ be the conjugates of
$x$, where $n=[\Q(x):\Q]\ge2$; each is taken by $d/n$ embeddings, so
$\Var(x)=n^{-2}\sum_{i<j}(x_i-x_j)^2$.  If $n=2$, then $(x_1-x_2)^2$ is the discriminant of a totally
real quadratic integer, so it is at least $5$.  If $n\ge3$, then the product of the $(x_i-x_j)^2$ is
a nonzero integer, and the inequality of arithmetic and geometric means gives
$\sum_{i<j}(x_i-x_j)^2\ge\binom n2$, so $\Var(x)\ge(n-1)/(2n)\ge1/3$.

Consequently the points of $\LamK$ are at distance at least $\sqrt{d/3}$ from each other, and a
packing argument gives $N(V)\le(1+2\sqrt{3V})^{d-1}$ for the number $N(V)$ of $x\in\OK/\Z$ with
$\Var(x)\le V$.  Integrating by parts,
$\sum_{x\ne0}e^{-a\Var(x)}=a\int_{1/3}^\infty e^{-aV}(N(V)-1)\,dV\le\varepsilon_{d,\ell}$.
\end{proof}

In the fields we computed, $\#T_1$ is close to $|\Delta_K|^{1/2}/(d-1)!$ (Section~\ref{sec:computations}),
so $c_{d,1}$ exceeds the truth by a factor of about $c_{d,1}(d-1)!$, that is, $7$ for $d=3$ and $28$
for $d=4$.

\begin{lemma}\label{lem:reg}
For every class $C\in\Cl$ we have $\zeta(2;C)\ge2^{d-2}R_K/(\gamma_d^d|\Delta_K|)$.
\end{lemma}

\begin{proof}
Fix a fractional ideal $\mathfrak f\in C^{-1}$.  The integral ideals in $C$ are the ideals
$\xi\mathfrak f^{-1}$ with $\xi\in(\mathfrak f\setminus0)/U$, so after scaling $\mathfrak f$ to a
lattice $\Lambda\subset\R^d$ of covolume $1$ we get $\zeta(2;C)=Z/|\Delta_K|$, where
$Z=\sum_{x\in(\Lambda\setminus0)/U}|\N(x)|^{-2}$.  We will show that $Z\ge2^{d-2}R_K/\gamma_d^d$ by
averaging over the orbit of $\Lambda$ under the diagonal group.

For $t\in H$ let $e^t$ act on $\R^d$ by $(e^{t_1},\dots,e^{t_d})$, and for $x$ with $\N(x)\ne0$ let
$\lambda(x)\in H$ have coordinates $\log|x_i|-\frac1d\log|\N(x)|$.  Then $\N(e^tx)=\N(x)$ and
$\lambda(e^tx)=\lambda(x)+t$.  Let $k(\lambda)=e^{-2dM(\lambda)}/Z_k$ on $H$, where
$M(\lambda)=\max_i\lambda_i\ge0$ and $Z_k=\int_He^{-2dM}$; the set $\{M\le s\}$ is a regular simplex of
volume $d^{d-1/2}s^{d-1}/(d-1)!$, so $Z_k=\sqrt d/2^{d-1}$.  Put $\varphi(x)=|\N(x)|^{-2}k(\lambda(x))$.
Then $\int_H\varphi(e^tx)\,dt=|\N(x)|^{-2}$.

For $u\in U^+$, multiplication by $u$ maps $\Lambda$ onto itself and equals $e^{\log u}$, where
$\log u=(\log\sigma_i(u))_i\in H$.  These vectors form a lattice $\log U^+$ in $H$ of covolume
$\sqrt d\,R_K[U:\pm U^+]$, and $U^+$ acts freely on $\Lambda\setminus0$.  For a fundamental domain $F$
of $\log U^+$, unfolding gives
\[
  \int_F\sum_{x\in\Lambda\setminus0}\varphi(e^tx)\,dt=\sum_{x\in(\Lambda\setminus0)/U^+}|\N(x)|^{-2}=[U:U^+]\,Z .
\]
For every $t$ the unimodular lattice $e^t\Lambda$ contains a nonzero $v$ with $|v|^2\le\gamma_d$.
Writing $|v_i|=|\N(v)|^{1/d}e^{\lambda_i}$ we get $|\N(v)|^{-2}\ge(\sum_ie^{2\lambda_i}/\gamma_d)^d\ge e^{2dM(\lambda)}/\gamma_d^d$,
hence $\varphi(v)\ge1/(\gamma_d^dZ_k)=2^{d-1}/(\sqrt d\,\gamma_d^d)$.  So the integrand is at least this
number, and $[U:U^+]=2[U:\pm U^+]$ gives $Z\ge2^{d-2}R_K/\gamma_d^d$.
\end{proof}

\section{Proof of Theorem \ref{thm:main}}\label{sec:proof}

By Proposition~\ref{prop:identity} and Lemma~\ref{lem:count},
\begin{equation}\label{eq:key}
  \zeta(2;[\Dif])=\frac{A_d\,\#T_{r_d}}{|\Delta_K|^{3/2}}\le\frac{A_dc_d}{|\Delta_K|} .
\end{equation}
An integral ideal $\mathfrak b\in[\Dif]$ contributes $\N\mathfrak b^{-2}$ to $\zeta(2;[\Dif])$, which
proves \ref{it:class}.  Lemma~\ref{lem:small} and \ref{it:class} give
$(|\Delta_K|/(A_dc_d))^{1/2}\le G_d|\Delta_K|^{(d-2)/(2d-2)}$, that is,
$|\Delta_K|^{1/(2d-2)}\le(A_dc_d)^{1/2}G_d$, which is \ref{it:disc}.  Lemma~\ref{lem:reg} with
$C=[\Dif]$ and \eqref{eq:key} give \ref{it:reg}.\qed

\medskip
Table~\ref{tab:bounds} lists the constants, computed with PARI/GP (Section~\ref{sec:computations}),
using the known values of $\gamma_n$ for $n\le8$ \cite{CS} and Blichfeldt's bound
$\gamma_n\le\frac2\pi\Gamma(2+\frac n2)^{2/n}$ \cite{Bl} for $n\ge9$.  For $d\ge6$ the table assumes
$m(K)=r_d$.

\begin{table}[ht]
\centering
\begin{tabular}{rrrrrrr}
\toprule
$d$ & $r_d$ & $A_d$ & $c_d$ & $G_d$ & bound \ref{it:disc} for $|\Delta_K|$ & bound \ref{it:reg} for $R_K$ \\
\midrule
2 & 1 & 6.494 & 2.073 & 1 & 13 & 18 \\
3 & 1 & 122.1 & 3.328 & 0.3849 & 3621 & 407 \\
4 & 1 & $5.061\cdot10^{3}$ & 4.597 & 0.25 & $3.1\cdot10^{9}$ & $2.3\cdot10^{4}$ \\
5 & 1 & $3.632\cdot10^{5}$ & 5.926 & 0.09499 & $1.4\cdot10^{17}$ & $2.2\cdot10^{6}$ \\
6 & 2 & $1.926\cdot10^{4}$ & 178.6 & 0.04724 & $2.7\cdot10^{19}$ & $4.6\cdot10^{6}$ \\
7 & 1 & $6.227\cdot10^{9}$ & 13.04 & 0.01770 & $2.7\cdot10^{44}$ & $1.6\cdot10^{11}$ \\
43 & 7 & $4.134\cdot10^{56}$ & $1.805\cdot10^{26}$ & $1.10\cdot10^{-17}$ & $10^{2056.2}$ & $3.6\cdot10^{105}$ \\
\bottomrule
\end{tabular}
\caption{Constants of Theorem \ref{thm:main}, to four significant digits; the discriminant bounds for
$d\le3$ are rounded down to integers, all other bounds are rounded up.
The file \file{results/etype/constants.txt} has all $d\le43$.}
\label{tab:bounds}
\end{table}

\begin{proof}[Proof of Corollary $\ref{cor:cubic}$]
For $d=2$ the bound is $|\Delta_K|\le13$.  The fields of discriminant $8$, $12$ and $13$ violate
\eqref{eq:condP}: in each we found $\nu\in T_1$ and a non-unit $\alpha\in\OK$ with
$\nu/\alpha\in\OKd$, and checked these three properties in exact arithmetic.  For $d=3$ the bound is
$|\Delta_K|\le3621$.  There are $122$ totally real cubic fields in this range (we compared the LMFDB
list \cite{LMFDB} with PARI's \file{nflist}).  For each of them except $\Q(\zeta_7+\zeta_7^{-1})$ we
found such a certificate $(\nu,\alpha)$, again checked exactly.  Finally, $\Q(\sqrt5)$ and
$\Q(\zeta_7+\zeta_7^{-1})$ satisfy \eqref{eq:condP} by Lemma~\ref{lem:etype}, since they admit
universal $\Z$-forms: $x^2+y^2+z^2$ over $\Q(\sqrt5)$ \cite{Ma} and $x^2+y^2+z^2+w^2+xy+xz+xw$ over
$\Q(\zeta_7+\zeta_7^{-1})$ \cite[Theorem~1.2]{KY1}.  The last sentence of the corollary follows from
Lemma~\ref{lem:etype}.
\end{proof}

\section{Computations}\label{sec:computations}

All computations used PARI/GP 2.15.4 and the scripts \file{src/etype\_zeta.gp},
\file{src/etype\_zeta\_checks.gp}, \file{src/etype\_zeta\_checks6.gp} and \file{src/etype\_census.sh};
the outputs are in \file{results/etype/}.  The sets $T_\ell$ were enumerated in floating point
(\file{qfminim} on the trace form of $\OKd$, with a margin of $10^{-9}$), and every element found was
then checked exactly.  Class groups come from \file{bnfinit} and are therefore conditional on GRH,
except where a statement is checked by an explicit certificate.

\medskip\noindent\emph{Proposition $\ref{prop:siegel}$.}  For the cubic fields $x^3-21x-28$,
$x^3-x^2-16x+8$ (both with $\Cl\cong\Z/3$ and non-principal $\Dif$) and $x^3-x^2-9x+8$ ($\Cl\cong\Z/3$,
principal $\Dif$), for the quartic field $x^4-x^3-19x^2+4x+76$ ($\Cl\cong\Z/2$), and for the sextic
field $x^6-x^5-6x^4+6x^3+8x^2-8x+1$ (class number $1$, $r_6=2$, so $T_2$ enters), the two sides of
\eqref{eq:siegelC} agree for every class to within $10^{-36}$.  The right side was computed from the
values $L(2,\chi)$ of the class group characters (\file{lfun}), the left side exactly.  For the
sextic field $x^6-x^5-18x^4+17x^3+58x^2-16x-1$ ($|\Delta_K|=46{,}411{,}625$, $\Cl\cong\Z/2$, principal
$\Dif$) they agree to within $10^{-32}$.

\medskip\noindent\emph{Condition \eqref{eq:condP} on field lists.}  We tested \eqref{eq:condP} with
$m(K)=1$ on the LMFDB lists of totally real fields \cite{LMFDB}: the first $40{,}000$ cubic fields
($|\Delta_K|\le741{,}448$), all $13{,}073$ quartic fields with $|\Delta_K|\le10^6$ and all $9{,}494$
quintic fields with $|\Delta_K|\le10^7$.  The fields satisfying \eqref{eq:condP} are exactly those of
discriminant $49$ (cubic), $725$ and $1125$ (quartic) and $14641$ (quintic), all with class number $1$.
Separately, among all $193{,}467$ cubic fields with $|\Delta_K|\le3{,}374{,}997$, the $9{,}859$
fields with non-principal $\Dif$ all violate \eqref{eq:condP}; on average $16.5\%$ of the elements
$\nu\in T_1$ have no principal divisor of $\nu\Dif$ other than $\OK$, but never all of them.  The same
holds for the $19$ quartic fields with non-principal $\Dif$; no quintic field in the list has one.

\medskip\noindent\emph{The estimates.}  In each of the ten fields of \file{results/etype/estimates.txt}, the element $\xi$ of
Lemma~\ref{lem:small} gives $\N(\xi\Dif)\le6$, far below $G_d|\Delta_K|^{(d-2)/(2d-2)}$; for instance
$\N(\xi\Dif)=6$ for $x^3-435x-2774$, where $|\Delta_K|=3{,}374{,}568$.  The numbers $\#T_1$ are close
to $|\Delta_K|^{1/2}/(d-1)!$ (for that field $\#T_1=919$ against $918.5$).  The smallest partial value
$\min_C\zeta(2;C)$ exceeded the bound of Lemma~\ref{lem:reg} by factors between $23$ and $9177$.  So
\ref{it:class} of Theorem~\ref{thm:main} fails by a wide margin in practice, and the size of the
bounds in Table~\ref{tab:bounds} comes from the uniform estimates, not from the method.

\section{Higher degrees}\label{sec:higher}

If $m(K)<r_d$, then \eqref{eq:siegelC} with $C=[\OK]$ reads
\begin{equation}\label{eq:general}
  \sum_{\ell=m(K)}^{r_d}b_\ell(2d)\,S_\ell=(-1)^d\frac{|\Delta_K|^{3/2}}{(4\pi^2)^d}\zeta(2;[\Dif]),
  \qquad S_\ell=\sum_{\nu\in T_\ell}\ \sum_{(\alpha)\mid\nu\Dif}\N(\alpha),
\end{equation}
and \eqref{eq:condP} gives $S_{m(K)}=\#T_{m(K)}$.  By Lemma~\ref{lem:small} the right side has
absolute value at least
\[
  \frac{|\Delta_K|^{1/2+1/(d-1)}}{(4\pi^2)^dG_d^2}.
\]
The signs of $(-1)^db_\ell(2d)$
for $d\le43$ (\file{results/etype/checks.txt}) are all positive for odd $d$ and for $d=6$; for even
$d\ge8$ they are negative for the smaller $\ell$ and positive for the last one or two.  Since
$S_{m(K)}=\#T_{m(K)}\le c_{d,m(K)}|\Delta_K|^{1/2}$, the argument of Section~\ref{sec:proof} gives an
explicit bound for $|\Delta_K|$ as soon as there are $\epsilon>0$ and $B$, depending only on $d$, with
\[
  \sum_{m(K)<\ell\le r_d,\ (-1)^db_\ell(2d)>0}|b_\ell(2d)|\,S_\ell\le B\,|\Delta_K|^{1/2+1/(d-1)-\epsilon}.
\]
Here $S_\ell$ counts pairs $(\delta,\alpha)$, up to $U^+$, with $\delta\in\OKd$ and $\alpha\in\OK$
totally positive, $\Tr(\delta\alpha)=\ell$, weighted by $\N(\alpha)$.  The minimizing property of
\cite[Lemma~4.6]{KY1} controls only the pairs with $\Tr(\delta\alpha)$ minimal for $\delta$.  When
$\OKd=\delta_0\OK$, the pairs with $\delta=\delta_0$ already give $S_2\ge\sum\N(\alpha)$ over the
totally positive $\alpha$ with $\Tr(\delta_0\alpha)=2$.  Heuristically there are about
$|\Delta_K|^{1/2}$ such $\alpha$, with norms up to $(2/d)^d|\Delta_K|$, so this sum is of order
$|\Delta_K|^{3/2}$.  The obstacle is therefore present also for principal codifferent, and it is the
reason for the hypothesis $m(K)=r_d$ in \cite[Theorem~2.4]{KM}.

What is needed is a description of the elements of $\OK$ in the shells
$\Tr(\delta\alpha)\le r_d\le8$ above the minimum.  For a universal $\Z$-form $Q$ these elements are
the values $Q(v)$ at the vectors $v$ of the tensor product of $Q$ with the trace form
$x\mapsto\Tr(\delta x^2)$ on $\OK$ whose norm is at most $r_d$ times the minimum.  Splitting of vectors of norm below twice
the minimum fails in general: in $A_2\otimes A_2$ the vector $x_1\otimes y_1+x_2\otimes y_2$ with
$\langle x_1,x_2\rangle=-1$, $\langle y_1,y_2\rangle=1$ has norm $6$, below $2\min(A_2\otimes A_2)=8$.
A replacement would have to use the universality of $Q$ (every element of the shell is represented)
and not only the E-type property.

\begin{thebibliography}{LMFDB}

\bibitem[Ar]{Ar}
J.~V. Armitage, On a theorem of Hecke in number fields and function fields,
Invent. Math. 2 (1967), 238--246.

\bibitem[Bl]{Bl}
H.~F. Blichfeldt, The minimum value of quadratic forms, and the closest packing of spheres,
Math. Ann. 101 (1929), 605--608.

\bibitem[CS]{CS}
J.~H. Conway, N.~J.~A. Sloane,
\emph{Sphere Packings, Lattices and Groups}, 3rd ed., Grundlehren Math. Wiss. 290,
Springer, 1999.

\bibitem[KL]{KL}
D.~Kim, S.~H. Lee, Lifting problem for universal quadratic forms over totally real cubic number
fields, arXiv:2307.07118.

\bibitem[KM]{KM}
V.~Kala, M.~Melistas, Universal quadratic forms and Dedekind zeta functions, arXiv:2311.12911.

\bibitem[KY1]{KY1}
V.~Kala, P.~Yatsyna, Lifting problem for universal quadratic forms,
Adv. Math. 377 (2021), 107497.

\bibitem[KY2]{KY2}
V.~Kala, P.~Yatsyna, Even better sums of squares over quintic and cyclotomic fields,
arXiv:2402.03850.

\bibitem[LMFDB]{LMFDB}
The LMFDB Collaboration, The L-functions and modular forms database, \url{https://www.lmfdb.org}.

\bibitem[Ma]{Ma}
H.~Maa\ss, \"Uber die Darstellung total positiver Zahlen des K\"orpers $R(\sqrt5)$ als Summe von drei
Quadraten, Abh. Math. Sem. Univ. Hamburg 14 (1941), 185--191.

\bibitem[Si]{Si}
C.~L. Siegel, Berechnung von Zetafunktionen an ganzzahligen Stellen,
Nachr. Akad. Wiss. G\"ottingen Math.-Phys. Kl. II (1969), 87--102.

\end{thebibliography}

\end{document}
